\section*{Erd\H{o}s Problem 1127: distinct distances and the continuum hypothesis}

\paragraph{Statement and answer.}
For each fixed positive integer \(n\), the assertion that \(\mathbb R^n\) is a countable union of sets with all unordered pairwise distances distinct is equivalent, over ZFC, to the continuum hypothesis (CH). We prove the obstruction and the one-dimensional construction below, and explicitly use Kunen's theorem for the higher-dimensional positive direction. This reconstructs the known independence result; see \url{https://www.erdosproblems.com/1127}.

\paragraph{CH gives a partition of the line.}
Assume CH and index a Hamel basis of \(\mathbb R\) over \(\mathbb Q\) as \(\{b_\alpha:\alpha<\omega_1\}\). For each countable ordinal \(\alpha\), fix an injection \(e_\alpha:\alpha\to\mathbb N\).
If \(x\ne0\), express it uniquely as a finite rational linear combination of basis elements, and let \(\alpha\) be the largest index appearing. Colour \(x\) by its nonzero leading coefficient together with the finite list of pairs
\[
 \{(e_\alpha(\beta),q_\beta):\beta<\alpha,\ q_\beta\ne0\}.
\]
There are countably many possible colours. In a fixed colour, at most one real has any given largest index \(\alpha\). Every finite set of distinct elements of that colour is therefore linearly independent: in a purported nontrivial rational relation, the uniquely largest basis index among the participating reals cannot cancel.
Give \(0\) a separate colour.

Rational linear independence implies distinct distances. An equality
\(|x-y|=|z-w|\) between two distinct unordered pairs, each consisting of distinct elements of one colour, yields \(x-y=\pm(z-w)\), a nontrivial rational linear relation after collecting repeated elements.

\paragraph{Failure of CH forbids such a partition, already on a line.}
If \(2^{\aleph_0}>\aleph_1\), a Hamel basis has cardinality \(2^{\aleph_0}\). Indeed the rational span of an infinite set has the same cardinality as that set. Choose disjoint subfamilies \(\{u_i:i\in I\}\), \(\{v_j:j\in J\}\) with \(|I|=\aleph_1\), \(|J|=\aleph_2\).
Given any countable colouring of the reals, for each \(j\) choose distinct \(i_j,i'_j\) such that \(u_{i_j}+v_j\) and \(u_{i'_j}+v_j\) have a common colour \(c_j\). Such a pair exists because \(I\) is uncountable.

There are only \(\aleph_1\) possible triples \((\{i_j,i'_j\},c_j)\). Two distinct \(j,j'\) therefore give the same triple. The four reals
\[
 u_i+v_j,\quad u_{i'}+v_j,\quad u_i+v_{j'},\quad u_{i'}+v_{j'}
\]
are distinct by basis independence, are monochromatic, and have two distinct pairs at distance \(|u_i-u_{i'}|\). This contradicts the required property.
A partition in any higher dimension would restrict to such a partition of a line, so the same obstruction applies.

\paragraph{Explicit higher-dimensional input.}
K. Kunen, \emph{Partitioning Euclidean space}, Mathematical Proceedings of the Cambridge Philosophical Society 102 (1987), 379--383, proves in Section 2: under CH, for every finite \(n\) there is a countable colouring of \(\mathbb R^n\) whose colour classes have distinct distances.
The primary paper is \url{https://doi.org/10.1017/S0305004100067426}.
That geometric transfinite-construction theorem is an external input here; the Hamel-basis argument above alone does not prove its higher-dimensional version. Together with the obstruction it proves the claimed equivalence for every positive \(n\).
The usual relative consistency and independence of CH are the additional foundational input for calling the question independent of ZFC.

\paragraph{Completion estimate.}
The equivalence is established, with Kunen's higher-dimensional theorem and the standard independence of CH explicitly imported.
\textbf{COMPLETION: 100\%}
